% !TEX program = xelatex
\documentclass[a4paper, 11pt]{article}

% ── Encoding & Fonts ──
\usepackage{fontspec}
\setmainfont{Noto Sans CJK KR}
\setsansfont{Noto Sans CJK KR}
\setmonofont{Noto Sans Mono CJK KR}[Scale=0.85, BoldFont=Noto Sans Mono CJK KR Bold]

% ── Page Layout ──
\usepackage[
  top=2.2cm, bottom=2.2cm, left=2.2cm, right=2.2cm,
  headheight=14pt
]{geometry}

% ── Colors (ink-saving: minimal, dark text only) ──
\usepackage{xcolor}
\definecolor{mainblue}{HTML}{1B3A5C}
\definecolor{codeborder}{HTML}{AAAAAA}
\definecolor{sectionbar}{HTML}{333333}
\definecolor{darktext}{HTML}{1F2937}
\definecolor{graytext}{HTML}{6B7280}
\definecolor{hintborder}{HTML}{888888}
\definecolor{typeA}{HTML}{2563EB}
\definecolor{typeB}{HTML}{059669}
\definecolor{typeC}{HTML}{DC2626}
\definecolor{typeD}{HTML}{92400E}

% ── Packages ──
\usepackage{fancyhdr}
\usepackage{enumitem}
\usepackage{listings}
\usepackage{tcolorbox}
\tcbuselibrary{skins,breakable}
\usepackage{tikz}
\usetikzlibrary{calc}
\usepackage{tabularx}
\usepackage{array}
\usepackage{setspace}
\usepackage{parskip}
\usepackage{lastpage}

% ── Paragraph & Spacing ──
\setlength{\parindent}{0pt}
\setlength{\parskip}{4pt}
\setstretch{1.25}

% ── Header / Footer ──
\pagestyle{fancy}
\fancyhf{}
\renewcommand{\headrulewidth}{0.4pt}
\fancyhead[L]{\small\color{graytext} 파이썬 요약 문제집 — 문제편}
\fancyhead[R]{\small\color{graytext} 고1 정보교과 대비}
\fancyfoot[C]{\small\color{graytext} \thepage\,/\,\pageref{LastPage}}

% ── Code Listing Style (white bg, thin border) ──
\lstdefinestyle{pycode}{
  language=Python,
  basicstyle=\ttfamily\small,
  keywordstyle=\bfseries\color{black},
  stringstyle=\color{darktext},
  commentstyle=\itshape\color{graytext},
  numberstyle=\tiny\color{graytext},
  numbers=left,
  numbersep=8pt,
  xleftmargin=14pt,
  framexleftmargin=14pt,
  backgroundcolor=\color{white},
  frame=single,
  rulecolor=\color{codeborder},
  breaklines=true,
  showstringspaces=false,
  tabsize=4,
  aboveskip=8pt,
  belowskip=8pt,
}
\lstset{style=pycode}

% ── Section Styling (outline only) ──
\newcommand{\unitsection}[2]{%
  \vspace{8pt}%
  \noindent%
  \begin{tikzpicture}
    \draw[sectionbar, line width=1.2pt] (0,0) rectangle (\textwidth, 0.95);
    \node[anchor=west, font=\large\bfseries\color{sectionbar}] at (0.4, 0.475) {#1};
    \node[anchor=east, font=\small\color{graytext}] at (\textwidth-0.3, 0.475) {#2};
  \end{tikzpicture}%
  \vspace{6pt}%
}

% ── Problem Box (white bg, thin border) ──
\newcounter{probcnt}
\newtcolorbox{problembox}[2][]{%
  enhanced,
  breakable,
  colback=white,
  colframe=codeborder,
  coltitle=black,
  fonttitle=\bfseries\small,
  title={\stepcounter{probcnt}\,문제 \theprobcnt\hfill #2},
  colbacktitle=white,
  left=8pt, right=8pt, top=6pt, bottom=6pt,
  boxrule=0.5pt,
  arc=2pt,
  toptitle=3pt, bottomtitle=3pt,
  titlerule=0.4pt,
  #1
}

% ── Type Tags (text only) ──
\newcommand{\tagpredict}{{\small\color{typeA}\bfseries [출력 예측]}}
\newcommand{\tagfill}{{\small\color{typeB}\bfseries [빈칸 채우기]}}
\newcommand{\tagerror}{{\small\color{typeC}\bfseries [오류 찾기]}}
\newcommand{\tagconcept}{{\small\color{typeD}\bfseries [개념]}}

% ── Hint Box (left rule only) ──
\newtcolorbox{hintbox}{%
  enhanced, breakable,
  colback=white,
  colframe=hintborder,
  boxrule=0pt, leftrule=2pt,
  left=8pt, right=8pt, top=4pt, bottom=4pt,
  arc=0pt,
  fontupper=\small\color{darktext},
}

% ── Answer blank ──
\newcommand{\answerblank}[1][4cm]{%
  \par\vspace{4pt}%
  \noindent 답:\,\rule[-.3ex]{#1}{0.5pt}%
  \vspace{2pt}%
}

\newcommand{\answerlines}[1][3]{%
  \par\vspace{4pt}%
  \foreach \i in {1,...,#1}{%
    \noindent\rule{\linewidth}{0.3pt}\par\vspace{6pt}%
  }%
}

% ══════════════════════════════════════════════════════════════
\begin{document}

% ── COVER PAGE (ink-saving: text only) ──
\thispagestyle{empty}

\vspace*{\fill}
\begin{center}
{\small\color{graytext} 고1 정보교과 대비}\\[20pt]
{\fontsize{36}{42}\selectfont\bfseries\color{mainblue} 파이썬}\\[8pt]
{\fontsize{24}{30}\selectfont\color{sectionbar} 요약 문제집}\\[20pt]
\rule{6cm}{0.8pt}\\[14pt]
{\fontsize{13}{18}\selectfont\color{graytext} 문제편}\\[30pt]
{\small\color{graytext} 총 37문제\enspace·\enspace 출력 예측 · 빈칸 채우기 · 오류 찾기 · 개념}
\end{center}
\vspace*{\fill}

\clearpage

% ── TABLE OF CONTENTS ──
\thispagestyle{fancy}
\vspace*{8pt}
{\Large\bfseries\color{mainblue} 목차}
\vspace{10pt}

\renewcommand{\arraystretch}{1.6}
\begin{tabularx}{\textwidth}{@{} >{\bfseries\color{mainblue}}l X >{\color{graytext}\small}r @{}}
1단원 & 입출력 기본 다지기 & 문제 1\,--\,6 \\
2단원 & 조건문 완전 정복 & 문제 7\,--\,13 \\
3단원 & 조건문 응용 프로젝트 & 문제 14\,--\,17 \\
4단원 & while 반복문과 누적 계산 & 문제 18\,--\,24 \\
5단원 & for 반복문과 패턴 출력 & 문제 25\,--\,30 \\
6단원 & 중첩 반복문 마스터 & 문제 31\,--\,34 \\
7단원 & 기초 모듈 활용 {\normalfont\color{graytext}(부록)} & 문제 35\,--\,37 \\
\end{tabularx}

\vspace{20pt}
\begin{hintbox}
\textbf{활용법}\\[2pt]
• 코드를 눈으로 읽고, 연필로 한 줄씩 추적하며 풀어보세요.\\
• 변수의 값 변화를 빈 공간에 메모하면 실수를 줄일 수 있습니다.\\
• 유형 태그(\,\tagpredict\;\tagfill\;\tagerror\;\tagconcept\,)를 참고해 약한 유형을 집중 연습하세요.
\end{hintbox}

\clearpage

% ══════════════════════════════════════════════════════════════
%  UNIT 1
% ══════════════════════════════════════════════════════════════
\unitsection{1단원\quad 입출력 기본 다지기}{6문제}

\begin{problembox}{\tagpredict}
다음 코드의 실행 결과를 쓰시오.
\begin{lstlisting}
a = "20"
b = "25"
print(a + b)
print(int(a) + int(b))
\end{lstlisting}
\answerlines[2]
\end{problembox}

\begin{problembox}{\tagpredict}
사용자가 \texttt{170823}을 입력했을 때, 다음 코드의 실행 결과를 쓰시오.
\begin{lstlisting}
data = input()
year = int(data[0:2]) + 2000
month = data[2:4]
day = data[4:6]
print(year, month, day, sep="/")
\end{lstlisting}
\answerblank[8cm]
\end{problembox}

\begin{problembox}{\tagfill}
아래 코드가 다음과 같이 출력되도록 \texttt{(가)}와 \texttt{(나)}에 들어갈 것을 쓰시오.

\vspace{2pt}
\textbf{[실행 결과]}\quad\texttt{사과\$바나나\$포도}

\begin{lstlisting}
fruits = ["사과", "바나나", "포도"]
print(fruits[0], fruits[1], fruits[2], sep=__(가)__, end=__(나)__)
\end{lstlisting}

(가):\,\rule[-.3ex]{4cm}{0.5pt}\qquad (나):\,\rule[-.3ex]{4cm}{0.5pt}
\end{problembox}

\begin{problembox}{\tagpredict}
사용자가 차례로 \texttt{7}과 \texttt{3}을 입력했을 때, 다음 코드의 실행 결과를 쓰시오.
\begin{lstlisting}
a = int(input())
b = int(input())
print(a // b, a % b)
print(a ** b)
\end{lstlisting}
\answerlines[2]
\end{problembox}

\begin{problembox}{\tagerror}
다음 코드를 실행하면 오류가 발생한다. 오류의 원인을 설명하고, 올바르게 수정하시오.
\begin{lstlisting}
age = input("나이를 입력하세요: ")
next_year = age + 1
print("내년 나이:", next_year)
\end{lstlisting}
\answerlines[3]
\end{problembox}

\begin{problembox}{\tagpredict}
사용자가 \texttt{4250}을 입력했을 때, 다음 코드의 실행 결과를 쓰시오.
\begin{lstlisting}
seconds = int(input())
hours = seconds // 3600
minutes = (seconds % 3600) // 60
secs = seconds % 60
print(hours, "시간", minutes, "분", secs, "초")
\end{lstlisting}
\answerblank[10cm]
\end{problembox}

% ══════════════════════════════════════════════════════════════
%  UNIT 2
% ══════════════════════════════════════════════════════════════
\clearpage
\unitsection{2단원\quad 조건문 완전 정복}{7문제}

\begin{problembox}{\tagconcept}
다음 두 코드의 실행 결과가 서로 다른 경우가 있다. \texttt{score = 72}일 때 각각의 결과를 쓰고, 차이가 생기는 이유를 간단히 설명하시오.

\begin{minipage}[t]{0.47\textwidth}
\textbf{[코드 A]}
\begin{lstlisting}
score = 72
if score >= 90:
    print("A")
elif score >= 80:
    print("B")
elif score >= 70:
    print("C")
else:
    print("D")
\end{lstlisting}
\end{minipage}
\hfill
\begin{minipage}[t]{0.47\textwidth}
\textbf{[코드 B]}
\begin{lstlisting}
score = 72
if score >= 90:
    print("A")
if score >= 80:
    print("B")
if score >= 70:
    print("C")
\end{lstlisting}
\end{minipage}

\vspace{6pt}
코드 A 결과:\,\rule[-.3ex]{3cm}{0.5pt}\qquad 코드 B 결과:\,\rule[-.3ex]{3cm}{0.5pt}

이유:\vspace{2pt}
\answerlines[2]
\end{problembox}

\begin{problembox}{\tagerror}
다음 코드에는 두 가지 오류가 있다. 모두 찾아 설명하시오.
\begin{lstlisting}
x = int(input())
if x > 0
    print("양수")
elif x = 0:
    print("영")
else:
    print("음수")
\end{lstlisting}
\answerlines[3]
\end{problembox}

\begin{problembox}{\tagpredict}
\texttt{x = 2024}일 때, 다음 코드의 실행 결과를 쓰시오.
\begin{lstlisting}
x = 2024
if x % 400 == 0:
    print("윤년")
elif x % 100 == 0:
    print("평년")
elif x % 4 == 0:
    print("윤년")
else:
    print("평년")
\end{lstlisting}
\answerblank[5cm]
\end{problembox}

\begin{problembox}{\tagfill}
세 변의 길이 \texttt{a}, \texttt{b}, \texttt{c}를 입력받아 삼각형의 종류를 판별하는 코드이다. 빈칸 \texttt{(가)}, \texttt{(나)}를 채우시오.
\begin{lstlisting}
a = int(input())
b = int(input())
c = int(input())
if a == b __(가)__ b == c:
    print("정삼각형")
elif a == b __(나)__ b == c or a == c:
    print("이등변삼각형")
else:
    print("일반 삼각형")
\end{lstlisting}

(가):\,\rule[-.3ex]{3cm}{0.5pt}\qquad (나):\,\rule[-.3ex]{3cm}{0.5pt}
\end{problembox}

\begin{problembox}{\tagpredict}
사용자가 \texttt{-15}를 입력했을 때, 다음 코드의 실행 결과를 쓰시오.
\begin{lstlisting}
n = int(input())
if n >= 0:
    if n % 2 == 0:
        print("양수 짝수")
    else:
        print("양수 홀수")
else:
    if n % 2 == 0:
        print("음수 짝수")
    else:
        print("음수 홀수")
\end{lstlisting}
\answerblank[5cm]
\end{problembox}

\begin{problembox}{\tagpredict}
다음 코드의 실행 결과를 쓰시오.
\begin{lstlisting}
a, b, c = 5, 3, 8
if not (a > b and b > c):
    print("조건1")
if a > c or b < c:
    print("조건2")
\end{lstlisting}
\answerlines[2]
\end{problembox}

\begin{problembox}{\tagfill}
좌표 \texttt{(x, y)}가 어느 사분면에 있는지 출력하는 코드이다. 빈칸을 채우시오.
\begin{lstlisting}
x = int(input())
y = int(input())
if x > 0 and y > 0:
    print("제1사분면")
elif __________ and y > 0:
    print("제2사분면")
elif x < 0 and __________:
    print("제3사분면")
else:
    print("제4사분면")
\end{lstlisting}
\answerlines[2]
\end{problembox}

% ══════════════════════════════════════════════════════════════
%  UNIT 3
% ══════════════════════════════════════════════════════════════
\clearpage
\unitsection{3단원\quad 조건문 응용 프로젝트}{4문제}

\begin{problembox}{\tagpredict}
사용자가 차례로 \texttt{170}과 \texttt{68}을 입력했을 때, 다음 코드의 실행 결과를 쓰시오.
\begin{lstlisting}
height = int(input()) / 100
weight = int(input())
bmi = weight / (height * height)
bmi = int(bmi * 10) / 10

if bmi >= 30:
    print(bmi, "고도비만")
elif bmi >= 25:
    print(bmi, "비만")
elif bmi >= 23:
    print(bmi, "과체중")
elif bmi >= 18.5:
    print(bmi, "정상")
else:
    print(bmi, "저체중")
\end{lstlisting}
\begin{hintbox}
\texttt{int(값 * 10) / 10}은 소수 첫째 자리까지만 남기는 테크닉입니다.
\end{hintbox}
\answerblank[8cm]
\end{problembox}

\begin{problembox}{\tagfill}
가위바위보 게임에서 승패를 판별하는 코드이다. 빈칸을 채우시오.\\
(\,\texttt{user}와 \texttt{com}은 각각 \texttt{"가위"}, \texttt{"바위"}, \texttt{"보"} 중 하나\,)
\begin{lstlisting}
if user == com:
    print("비김")
elif (user == "가위" and com == __________) or \
     (user == "바위" and com == __________) or \
     (user == "보" and com == __________):
    print("승리")
else:
    print("패배")
\end{lstlisting}
\answerlines[2]
\end{problembox}

\begin{problembox}{\tagpredict}
사용자가 차례로 \texttt{10}, \texttt{+}, \texttt{3}을 입력했을 때, 다음 코드의 실행 결과를 쓰시오.\\
그 다음, 사용자가 \texttt{10}, \texttt{\%}, \texttt{3}을 입력했을 때의 결과도 쓰시오.
\begin{lstlisting}
a = int(input())
op = input()
b = int(input())
if op == "+":
    print(a + b)
elif op == "-":
    print(a - b)
elif op == "*":
    print(a * b)
elif op == "/":
    print(a / b)
else:
    print("지원하지 않는 연산입니다")
\end{lstlisting}
입력 \texttt{10, +, 3}\;\;결과:\,\rule[-.3ex]{4cm}{0.5pt}

입력 \texttt{10, \%, 3}\;\;결과:\,\rule[-.3ex]{4cm}{0.5pt}
\end{problembox}

\begin{problembox}{\tagconcept}
다음 코드에서 \texttt{money = 3500}일 때 실행 결과를 쓰시오. 또한, \texttt{elif}를 모두 \texttt{if}로 바꾸면 결과가 어떻게 달라지는지 설명하시오.
\begin{lstlisting}
money = 3500
if money >= 5000:
    print("스테이크")
elif money >= 3000:
    print("김치찌개")
elif money >= 1500:
    print("라면")
else:
    print("물")
\end{lstlisting}
\answerlines[3]
\end{problembox}

% ══════════════════════════════════════════════════════════════
%  UNIT 4
% ══════════════════════════════════════════════════════════════
\clearpage
\unitsection{4단원\quad while 반복문과 누적 계산}{7문제}

\begin{problembox}{\tagpredict}
다음 코드의 실행 결과를 쓰시오.
\begin{lstlisting}
n = 1234
total = 0
while n > 0:
    total = total + n % 10
    n = n // 10
print(total)
\end{lstlisting}
\begin{hintbox}
변수 \texttt{n}과 \texttt{total}의 값 변화를 단계별로 추적해 보세요.
\end{hintbox}
\answerblank[4cm]
\end{problembox}

\begin{problembox}{\tagpredict}
다음 코드의 실행 결과를 쓰시오.
\begin{lstlisting}
i = 1
result = 1
while i <= 5:
    result = result * i
    i = i + 1
print(result)
\end{lstlisting}
\answerblank[4cm]
\end{problembox}

\begin{problembox}{\tagfill}
사용자가 \texttt{-1}을 입력할 때까지 숫자를 계속 입력받아 합계를 구하는 코드이다. 빈칸 \texttt{(가)}, \texttt{(나)}를 채우시오.
\begin{lstlisting}
total = 0
num = int(input())
while __(가)__:
    total = total + num
    num = __(나)__
print("합계:", total)
\end{lstlisting}
(가):\,\rule[-.3ex]{4cm}{0.5pt}\qquad (나):\,\rule[-.3ex]{4cm}{0.5pt}
\end{problembox}

\begin{problembox}{\tagpredict}
다음 코드의 실행 결과를 쓰시오.
\begin{lstlisting}
n = 36
i = 1
count = 0
while i <= n:
    if n % i == 0:
        count = count + 1
    i = i + 1
print(count)
\end{lstlisting}
\answerblank[4cm]
\end{problembox}

\begin{problembox}{\tagpredict}
다음 코드의 실행 결과를 쓰시오.
\begin{lstlisting}
n = 50
i = 1
total = 0
while i <= n:
    if i % 10 == 0:
        total = total + i
    i = i + 1
print(total)
\end{lstlisting}
\answerblank[4cm]
\end{problembox}

\begin{problembox}{\tagpredict}
다음 코드의 실행 결과를 쓰시오.
\begin{lstlisting}
n = 7
i = 2
is_prime = True
while i < n:
    if n % i == 0:
        is_prime = False
    i = i + 1
if is_prime:
    print("소수")
else:
    print("소수 아님")
\end{lstlisting}
\answerblank[4cm]
\end{problembox}

\begin{problembox}{\tagfill}
1부터 \texttt{n}까지의 수 중 3의 배수이면서 5의 배수가 아닌 수의 합을 구하는 코드이다. 빈칸을 채우시오.
\begin{lstlisting}
n = 30
i = 1
total = 0
while i <= n:
    if __________ and __________:
        total = total + i
    i = i + 1
print(total)
\end{lstlisting}
\answerlines[2]
\end{problembox}

% ══════════════════════════════════════════════════════════════
%  UNIT 5
% ══════════════════════════════════════════════════════════════
\clearpage
\unitsection{5단원\quad for 반복문과 패턴 출력}{6문제}

\begin{problembox}{\tagconcept}
다음 세 \texttt{range}가 생성하는 수열을 각각 나열하시오.

\begin{tabularx}{\textwidth}{@{} l X @{}}
(1) \texttt{range(1, 10, 2)} & 답:\,\rule[-.3ex]{6cm}{0.5pt} \\[6pt]
(2) \texttt{range(10, 0, -3)} & 답:\,\rule[-.3ex]{6cm}{0.5pt} \\[6pt]
(3) \texttt{range(5)} & 답:\,\rule[-.3ex]{6cm}{0.5pt} \\
\end{tabularx}
\end{problembox}

\begin{problembox}{\tagpredict}
다음 코드의 실행 결과를 쓰시오.
\begin{lstlisting}
total = 0
for i in range(1, 101):
    if i % 3 == 0 and i % 5 == 0:
        total = total + i
print(total)
\end{lstlisting}
\answerblank[4cm]
\end{problembox}

\begin{problembox}{\tagpredict}
\texttt{n = 24}일 때, 다음 코드의 실행 결과를 쓰시오.
\begin{lstlisting}
n = 24
gcd = 1
for i in range(1, min(n, 18) + 1):
    if n % i == 0 and 18 % i == 0:
        gcd = i
print(gcd)
\end{lstlisting}
\answerblank[4cm]
\end{problembox}

\begin{problembox}{\tagfill}
입력된 수 \texttt{n}의 약수를 모두 출력하되, 한 줄에 공백으로 구분하여 출력하는 코드이다. 빈칸을 채우시오.
\begin{lstlisting}
n = int(input())
for i in range(1, __________):
    if n % i == 0:
        print(i, end=__________)
\end{lstlisting}
\answerlines[2]
\end{problembox}

\begin{problembox}{\tagpredict}
다음 코드의 실행 결과를 쓰시오.
\begin{lstlisting}
for i in range(1, 6):
    if i == 3:
        continue
    if i == 5:
        break
    print(i, end=" ")
\end{lstlisting}
\answerblank[6cm]
\end{problembox}

\begin{problembox}{\tagpredict}
다음 코드의 실행 결과를 쓰시오.
\begin{lstlisting}
for i in range(2, 6):
    count = 0
    for j in range(1, i + 1):
        if i % j == 0:
            count = count + 1
    if count == 2:
        print(i, end=" ")
\end{lstlisting}
\answerblank[6cm]
\end{problembox}

% ══════════════════════════════════════════════════════════════
%  UNIT 6
% ══════════════════════════════════════════════════════════════
\clearpage
\unitsection{6단원\quad 중첩 반복문 마스터}{4문제}

\begin{problembox}{\tagpredict}
다음 코드의 실행 결과를 쓰시오.
\begin{lstlisting}
for i in range(1, 5):
    for j in range(1, 5):
        if i == j:
            print(1, end=" ")
        else:
            print(0, end=" ")
    print()
\end{lstlisting}
\answerlines[4]
\end{problembox}

\begin{problembox}{\tagpredict}
다음 코드의 실행 결과를 쓰시오.
\begin{lstlisting}
for i in range(1, 6):
    for j in range(i, 6):
        print("*", end="")
    print()
\end{lstlisting}
\answerlines[5]
\end{problembox}

\begin{problembox}{\tagfill}
다음 출력을 만들어내는 코드의 빈칸을 채우시오.

\textbf{[출력]}\\
\texttt{1}\\
\texttt{1 2}\\
\texttt{1 2 3}\\
\texttt{1 2 3 4}

\begin{lstlisting}
for i in range(1, 5):
    for j in range(1, __________):
        print(j, end=" ")
    print()
\end{lstlisting}
\answerblank[6cm]
\end{problembox}

\begin{problembox}{\tagpredict}
다음 코드의 실행 결과를 쓰시오.
\begin{lstlisting}
for i in range(2, 10, 2):
    for j in range(1, 4):
        print(i, "x", j, "=", i * j)
    print("---")
\end{lstlisting}
\begin{hintbox}
\texttt{range(2, 10, 2)}는 짝수 단만 생성합니다. 안쪽 반복은 1\,--\,3까지만.
\end{hintbox}
\answerlines[6]
\end{problembox}

% ══════════════════════════════════════════════════════════════
%  UNIT 7
% ══════════════════════════════════════════════════════════════
\clearpage
\unitsection{7단원\quad 기초 모듈 활용\;{\normalsize\color{graytext}(부록)}}{3문제}

\begin{problembox}{\tagconcept}
다음 코드의 실행 결과를 쓰시오.
\begin{lstlisting}
import math
print(math.ceil(3.2))
print(math.floor(3.8))
print(math.floor(-2.3))
\end{lstlisting}
\answerlines[3]
\end{problembox}

\begin{problembox}{\tagfill}
\texttt{random} 모듈을 사용하여 1부터 6까지의 정수를 무작위로 하나 출력하는 코드이다. 빈칸을 채우시오.
\begin{lstlisting}
import random
dice = random.__(가)__(__(나)__, __(다)__)
print("주사위:", dice)
\end{lstlisting}
(가):\,\rule[-.3ex]{3cm}{0.5pt}\quad (나):\,\rule[-.3ex]{2cm}{0.5pt}\quad (다):\,\rule[-.3ex]{2cm}{0.5pt}
\end{problembox}

\begin{problembox}{\tagpredict}
다음 코드의 실행 결과를 쓰시오.
\begin{lstlisting}
import math
r = 5
area = math.pi * r ** 2
area = round(area, 2)
print(area)
\end{lstlisting}
\begin{hintbox}
\texttt{math.pi}는 3.141592653589793…입니다. \texttt{round(값, 2)}는 소수 둘째 자리 반올림.
\end{hintbox}
\answerblank[5cm]
\end{problembox}

\vfill
\begin{center}
\color{graytext}\small — 문제편 끝 —
\end{center}

\end{document}
